\documentclass[11pt,a4paper]{article}

%=========================================================
% PACKAGES
%=========================================================
\usepackage[utf8]{inputenc}
\usepackage[T1]{fontenc}
\usepackage[french]{babel}
\usepackage{lmodern}
\usepackage{geometry}
\geometry{margin=2.1cm}
\usepackage{amsmath,amssymb,amsfonts}
\usepackage{mathtools}
\usepackage{physics}
\usepackage{siunitx}
\sisetup{locale = FR, per-mode = symbol}
\usepackage{graphicx}
\usepackage{xcolor}
\usepackage[most]{tcolorbox}
\usepackage{enumitem}
\usepackage{hyperref}
\usepackage{fancyhdr}
\usepackage{titlesec}
\usepackage{array}
\usepackage{booktabs}
\usepackage{multicol}

%=========================================================
% STYLE
%=========================================================
\hypersetup{
    colorlinks=true,
    linkcolor=blue!60!black,
    urlcolor=blue!60!black
}

\pagestyle{fancy}
\setlength{\headheight}{14pt}
\fancyhf{}
\lhead{\textsc{Grand Oral -- Balle de tennis}}
\rhead{\textsc{Approfondissement physique}}
\cfoot{\thepage}

\titleformat{\section}{\Large\bfseries\color{blue!55!black}}{}{0em}{}
\titleformat{\subsection}{\large\bfseries\color{blue!40!black}}{}{0em}{}
\titleformat{\subsubsection}{\bfseries\color{blue!30!black}}{}{0em}{}

\tcbset{
    arc=2mm,
    boxrule=0.7pt,
    breakable
}

\newtcolorbox{idee}[1]{
    title=\textbf{Idée physique : #1},
    colback=green!5,
    colframe=green!45!black
}

\newtcolorbox{modele}[1]{
    title=\textbf{Modèle : #1},
    colback=blue!4,
    colframe=blue!55!black
}

\newtcolorbox{bilan}[1]{
    title=\textbf{Bilan des forces : #1},
    colback=purple!4,
    colframe=purple!60!black
}

\newtcolorbox{equationbox}[1]{
    title=\textbf{Mise en équation : #1},
    colback=yellow!6,
    colframe=orange!80!black
}

\newtcolorbox{tennis}[1]{
    title=\textbf{Conclusion tennis : #1},
    colback=red!4,
    colframe=red!60!black
}

\newtcolorbox{oral}[1]{
    title=\textbf{Rédaction possible à l'oral : #1},
    colback=gray!7,
    colframe=black!60
}

\newtcolorbox{prepa}[1]{
    title=\textbf{Ouverture prépa accessible : #1},
    colback=cyan!4,
    colframe=cyan!55!black
}

%=========================================================
% DOCUMENT
%=========================================================
\begin{document}

\begin{center}
    {\Huge \textbf{Physique d'une balle de tennis}}\\[0.2cm]
    {\Large Trajectoire, lift, rebond, vent et cordage}\\[0.2cm]
    {\large Dossier rigoureux pour Grand Oral, niveau Terminale avec ouvertures vers la prépa}\\[0.4cm]
    \rule{0.85\linewidth}{0.5pt}
\end{center}

\vspace{0.4cm}

\tableofcontents

\newpage

%=========================================================
\section{Introduction générale}

Le tennis paraît être un sport de vitesse et de précision. Pourtant, derrière un simple échange, une balle de tennis met en jeu de nombreux phénomènes physiques : poids, frottements, rotation, effet Magnus, rebond, énergie, quantité de mouvement, vent, interaction avec les cordes de la raquette.

\medskip

Une balle de tennis est un objet particulièrement intéressant car elle est :
\begin{itemize}
    \item légère, de masse environ $m \approx \SI{58}{g}$ ;
    \item rapide, pouvant dépasser $\SI{50}{m.s^{-1}}$ au service ;
    \item rugueuse, grâce au feutre ;
    \item déformable ;
    \item souvent en rotation très rapide.
\end{itemize}

\begin{center}
\Large
\textbf{Problématique possible :}\\[0.2cm]
\textit{Comment la physique permet-elle de comprendre et d'optimiser la trajectoire d'une balle de tennis ?}
\end{center}

\medskip

L'objectif de ce dossier n'est pas seulement de décrire qualitativement les phénomènes. On veut, pour chaque piste, poser un modèle physique, écrire les forces, obtenir une équation, puis interpréter le résultat pour le tennis.

%=========================================================
\section{Méthode générale de modélisation}

Pour étudier une balle de tennis, on applique toujours la même démarche :

\begin{enumerate}
    \item \textbf{Définir le système étudié.}
    
    Ici, le système est souvent la balle de tennis assimilée à un point matériel ou à une sphère en rotation.
    
    \item \textbf{Choisir un référentiel.}
    
    On choisit en général le référentiel terrestre, supposé galiléen à l'échelle du mouvement.
    
    \item \textbf{Faire le bilan des forces.}
    
    On liste les forces extérieures : poids, frottements de l'air, force de Magnus, réaction du sol, force de frottement du sol, force de la raquette.
    
    \item \textbf{Appliquer la deuxième loi de Newton.}
    
    \[
    \sum \vec{F}_{ext}=m\vec{a}.
    \]
    
    \item \textbf{Résoudre ou exploiter les équations obtenues.}
    
    On peut obtenir une position, une vitesse, une accélération ou une relation énergétique.
    
    \item \textbf{Interpréter le résultat pour le tennis.}
\end{enumerate}

\begin{oral}{phrase utile}
L'intérêt de la modélisation est de remplacer une situation réelle complexe par un modèle simplifié, mais suffisamment précis pour expliquer les phénomènes principaux. En tennis, cela permet de comprendre pourquoi une balle liftée retombe plus vite, pourquoi le vent gêne le joueur, ou encore pourquoi la tension du cordage influence la puissance et le contrôle.
\end{oral}

%=========================================================
\section{Piste 1 -- Trajectoire sans frottements : la balle comme projectile}

\subsection{Idée physique}

\begin{idee}{premier modèle}
On commence par le modèle le plus simple : la balle est assimilée à un point matériel soumis uniquement à son poids. Ce modèle est celui de la chute libre avec vitesse initiale, étudié en Terminale.
\end{idee}

\subsection{Modèle}

\begin{modele}{balle sans frottements}
On considère une balle de masse $m$, lancée depuis une hauteur $y_0$ avec une vitesse initiale $\vec{v}_0$ faisant un angle $\alpha$ avec l'horizontale.

On choisit un repère $(O,\vec{i},\vec{j})$ :
\begin{itemize}
    \item $\vec{i}$ horizontal, dans le sens du mouvement ;
    \item $\vec{j}$ vertical vers le haut.
\end{itemize}

La vitesse initiale s'écrit :
\[
\vec{v}_0=v_0\cos(\alpha)\vec{i}+v_0\sin(\alpha)\vec{j}.
\]
\end{modele}

\subsection{Bilan des forces}

\begin{bilan}{poids seul}
On néglige les frottements de l'air. La seule force extérieure est le poids :

\[
\vec{P}=m\vec{g}=-mg\vec{j}.
\]
\end{bilan}

\subsection{Mise en équation}

\begin{equationbox}{deuxième loi de Newton}
D'après la deuxième loi de Newton :

\[
m\vec{a}=\vec{P}.
\]

Donc :

\[
m\vec{a}=-mg\vec{j}.
\]

En simplifiant par $m$ :

\[
\vec{a}=-g\vec{j}.
\]

Ainsi :

\[
a_x=0,
\qquad
a_y=-g.
\]
\end{equationbox}

\subsection{Résolution}

On intègre les composantes de l'accélération.

Pour la vitesse horizontale :

\[
a_x=0
\Rightarrow
v_x(t)=v_0\cos(\alpha).
\]

Pour la vitesse verticale :

\[
a_y=-g
\Rightarrow
v_y(t)=v_0\sin(\alpha)-gt.
\]

On intègre ensuite pour obtenir la position :

\[
x(t)=v_0\cos(\alpha)t,
\]

\[
y(t)=y_0+v_0\sin(\alpha)t-\frac12gt^2.
\]

On élimine $t$ grâce à :

\[
t=\frac{x}{v_0\cos(\alpha)}.
\]

Donc :

\[
y(x)=y_0+x\tan(\alpha)-\frac{g}{2v_0^2\cos^2(\alpha)}x^2.
\]

La trajectoire est une parabole.

\subsection{Exploitation physique}

La hauteur maximale est atteinte lorsque :

\[
v_y(t)=0.
\]

Donc :

\[
v_0\sin(\alpha)-gt=0,
\]

d'où :

\[
t_{\max}=\frac{v_0\sin(\alpha)}{g}.
\]

La hauteur maximale vaut alors :

\[
y_{\max}
=
y_0+\frac{v_0^2\sin^2(\alpha)}{2g}.
\]

\subsection{Conclusion tennis}

\begin{tennis}{utilité et limite du modèle}
Ce modèle montre que la balle retombe naturellement à cause du poids. Il permet de comprendre l'influence de la vitesse initiale et de l'angle de frappe.

Mais il est insuffisant pour le tennis réel :
\begin{itemize}
    \item il ne prend pas en compte les frottements de l'air ;
    \item il ne prend pas en compte la rotation ;
    \item il ne permet pas d'expliquer le lift ;
    \item il prédit une trajectoire trop idéale.
\end{itemize}

Au tennis, une balle frappée très fort et sans effet aurait tendance à sortir. Les joueurs doivent donc utiliser d'autres phénomènes physiques, notamment les frottements et l'effet Magnus.
\end{tennis}

\begin{oral}{rédaction possible}
Dans un premier modèle, on assimile la balle à un point matériel soumis uniquement à son poids. La deuxième loi de Newton conduit à une accélération verticale constante égale à $-g$. En intégrant, on obtient les équations horaires de la position, puis une trajectoire parabolique. Ce modèle est utile pour comprendre la base du mouvement, mais il ne suffit pas pour expliquer le tennis réel, car une balle rapide subit fortement les frottements de l'air et peut tourner sur elle-même.
\end{oral}

%=========================================================
\section{Piste 2 -- Frottements de l'air : pourquoi la balle ralentit-elle ?}

\subsection{Idée physique}

\begin{idee}{traînée aérodynamique}
Une balle de tennis traverse l'air. L'air exerce une force opposée au mouvement, appelée force de traînée. Pour une balle rapide, cette force est approximativement proportionnelle au carré de la vitesse.
\end{idee}

\subsection{Modèle}

\begin{modele}{mouvement horizontal simplifié}
Pour obtenir une équation exploitable à l'oral, on commence par un mouvement horizontal.

On suppose que la balle se déplace selon l'axe $x$ avec une vitesse $v(t)>0$.

On néglige ici la pesanteur pour étudier uniquement le ralentissement horizontal.
\end{modele}

\subsection{Bilan des forces}

\begin{bilan}{traînée quadratique}
La force de frottement de l'air est :

\[
\vec{F}_D=-kv^2\vec{i},
\]

où :

\[
k=\frac12\rho C_D S.
\]

Avec :
\begin{itemize}
    \item $\rho$ : masse volumique de l'air ;
    \item $C_D$ : coefficient de traînée ;
    \item $S$ : surface frontale de la balle.
\end{itemize}
\end{bilan}

\subsection{Mise en équation}

\begin{equationbox}{ralentissement horizontal}
La deuxième loi de Newton donne :

\[
m\frac{dv}{dt}=-kv^2.
\]

Donc :

\[
\frac{dv}{dt}=-\frac{k}{m}v^2.
\]
\end{equationbox}

\subsection{Résolution accessible}

On sépare les variables :

\[
\frac{dv}{v^2}=-\frac{k}{m}dt.
\]

Or :

\[
\int \frac{dv}{v^2}=\int v^{-2}dv=-\frac{1}{v}.
\]

Donc :

\[
-\frac1v=-\frac{k}{m}t+C.
\]

À $t=0$, $v(0)=v_0$, donc :

\[
C=-\frac1{v_0}.
\]

Ainsi :

\[
-\frac1v=-\frac{k}{m}t-\frac1{v_0}.
\]

En multipliant par $-1$ :

\[
\frac1v=\frac{k}{m}t+\frac1{v_0}.
\]

Donc :

\[
v(t)=\frac{v_0}{1+\frac{k v_0}{m}t}.
\]

\subsection{Position}

Comme :

\[
v(t)=\frac{dx}{dt},
\]

on a :

\[
x(t)=\int_0^t \frac{v_0}{1+\frac{k v_0}{m}u}\,du.
\]

En posant :

\[
a=\frac{k v_0}{m},
\]

on obtient :

\[
x(t)=\frac{v_0}{a}\ln(1+at).
\]

Donc :

\[
x(t)=\frac{m}{k}\ln\left(1+\frac{k v_0}{m}t\right).
\]

\subsection{Interprétation}

Sans frottement :

\[
x(t)=v_0t.
\]

Avec frottement quadratique :

\[
x(t)=\frac{m}{k}\ln\left(1+\frac{k v_0}{m}t\right).
\]

La distance augmente moins vite qu'en l'absence de frottement.

\subsection{Conclusion tennis}

\begin{tennis}{impact de la traînée}
La balle ralentit fortement pendant son vol. Plus elle est rapide au départ, plus la force de traînée est grande, car :

\[
F_D=kv^2.
\]

Donc doubler la vitesse multiplie environ la force de frottement par quatre.

Cela explique pourquoi :
\begin{itemize}
    \item un service très rapide perd beaucoup de vitesse avant d'arriver au relanceur ;
    \item une balle liftée ou très frappée ne conserve pas toute son énergie ;
    \item la précision du modèle parabolique est limitée au tennis.
\end{itemize}
\end{tennis}

\begin{oral}{rédaction possible}
Si l'on ajoute les frottements de l'air, la balle ne conserve plus une vitesse horizontale constante. Pour une balle rapide, on peut modéliser la traînée par une force proportionnelle à $v^2$ et opposée au mouvement. En résolvant l'équation différentielle obtenue, on trouve que la vitesse décroît comme une fonction du type $\frac{v_0}{1+at}$. La balle ralentit donc rapidement, ce qui est essentiel au tennis, notamment pour les services puissants.
\end{oral}

%=========================================================
\section{Piste 3 -- Effet Magnus et lift : pourquoi la balle plonge-t-elle ?}

\subsection{Idée physique}

\begin{idee}{rotation et force de Magnus}
Lorsqu'une balle tourne sur elle-même en avançant, l'écoulement de l'air autour d'elle devient dissymétrique. Il apparaît alors une force perpendiculaire à la vitesse : la force de Magnus. Pour une balle liftée, cette force est dirigée vers le bas.
\end{idee}

\subsection{Modèle}

\begin{modele}{lift simplifié}
On étudie une balle avançant horizontalement à vitesse approximativement constante $v_x$.

La balle tourne vers l'avant avec une vitesse angulaire $\omega$.

On suppose que la force de Magnus verticale est proportionnelle à $v^2$ :

\[
F_M = k_M v^2,
\]

avec :

\[
k_M=\frac12\rho C_L S.
\]

Dans le cas du lift, la force est dirigée vers le bas.
\end{modele}

\subsection{Bilan des forces}

\begin{bilan}{poids et Magnus}
Les forces verticales sont :
\begin{itemize}
    \item le poids :
    \[
    \vec{P}=-mg\vec{j};
    \]
    \item la force de Magnus, dirigée vers le bas pour le lift :
    \[
    \vec{F}_M=-k_Mv^2\vec{j}.
    \]
\end{itemize}

On néglige ici la traînée verticale pour garder un modèle accessible.
\end{bilan}

\subsection{Mise en équation}

\begin{equationbox}{accélération verticale effective}
D'après la deuxième loi de Newton :

\[
m a_y=-mg-k_Mv^2.
\]

Donc :

\[
a_y=-g-\frac{k_M}{m}v^2.
\]

On pose :

\[
g_{\text{eff}}=g+\frac{k_M}{m}v^2.
\]

Alors :

\[
a_y=-g_{\text{eff}}.
\]
\end{equationbox}

\subsection{Résolution}

Si $v$ est approximativement constante pendant une courte durée, alors $g_{\text{eff}}$ est constante.

On obtient donc :

\[
v_y(t)=v_{y0}-g_{\text{eff}}t,
\]

puis :

\[
y(t)=y_0+v_{y0}t-\frac12g_{\text{eff}}t^2.
\]

Comparé au cas sans lift :

\[
y_{\text{sans lift}}(t)=y_0+v_{y0}t-\frac12gt^2.
\]

Comme :

\[
g_{\text{eff}}>g,
\]

on a :

\[
y_{\text{lift}}(t)<y_{\text{sans lift}}(t)
\]

pour un même instant $t>0$.

\subsection{Conclusion tennis}

\begin{tennis}{effet du lift}
Le lift augmente l'accélération vers le bas. La balle retombe donc plus vite qu'une balle frappée à plat.

Cela permet au joueur :
\begin{itemize}
    \item de frapper plus fort sans sortir la balle ;
    \item de passer plus haut au-dessus du filet ;
    \item de créer une trajectoire bombée ;
    \item d'obtenir un rebond plus haut et plus agressif.
\end{itemize}

Le lift est donc une stratégie physique de contrôle.
\end{tennis}

\begin{oral}{rédaction possible}
Pour une balle liftée, la rotation crée une force de Magnus dirigée vers le bas. L'équation verticale devient alors $ma_y=-mg-F_M$. On peut interpréter cela comme une augmentation de la pesanteur apparente : la balle subit une accélération verticale plus grande que $g$. Elle retombe donc plus vite qu'une balle sans effet, ce qui permet au joueur de frapper fort tout en gardant la balle dans le court.
\end{oral}

\begin{prepa}{forme vectorielle}
En prépa, la direction de la force de Magnus peut être reliée au produit vectoriel :

\[
\vec{F}_M \propto \vec{\omega}\times \vec{v}.
\]

Cette écriture montre que la force dépend à la fois de l'axe de rotation et de la direction de la vitesse.
\end{prepa}

%=========================================================
\section{Piste 4 -- Slice et balle coupée : pourquoi la balle flotte-t-elle ?}

\subsection{Idée physique}

\begin{idee}{rotation arrière}
Dans un slice, la balle tourne souvent vers l'arrière. Le sens de la force de Magnus change. La balle peut alors avoir une trajectoire plus flottante, être davantage freinée et rebondir plus bas.
\end{idee}

\subsection{Modèle}

\begin{modele}{balle slicée}
On considère une balle avançant horizontalement avec une rotation arrière.

On modélise la force de Magnus verticale par :

\[
\vec{F}_M=+k_Mv^2\vec{j}.
\]

Ici, contrairement au lift, la force de Magnus est orientée vers le haut.
\end{modele}

\subsection{Bilan des forces}

\begin{bilan}{poids et Magnus vers le haut}
Les forces verticales sont :
\[
\vec{P}=-mg\vec{j},
\]

et :

\[
\vec{F}_M=+k_Mv^2\vec{j}.
\]

Donc :

\[
\sum \vec{F}=(-mg+k_Mv^2)\vec{j}.
\]
\end{bilan}

\subsection{Mise en équation}

\begin{equationbox}{accélération verticale réduite}
On applique la deuxième loi de Newton :

\[
ma_y=-mg+k_Mv^2.
\]

Donc :

\[
a_y=-g+\frac{k_M}{m}v^2.
\]

On peut écrire :

\[
a_y=-g_{\text{eff}},
\]

avec :

\[
g_{\text{eff}}=g-\frac{k_M}{m}v^2.
\]

Si la force de Magnus vers le haut est importante, $g_{\text{eff}}$ diminue.
\end{equationbox}

\subsection{Résolution}

On obtient :

\[
v_y(t)=v_{y0}-g_{\text{eff}}t,
\]

et :

\[
y(t)=y_0+v_{y0}t-\frac12g_{\text{eff}}t^2.
\]

Comme :

\[
g_{\text{eff}}<g,
\]

la balle tombe moins vite qu'une balle sans effet, dans ce modèle.

\subsection{Conclusion tennis}

\begin{tennis}{impact du slice}
Un slice peut donner l'impression que la balle flotte davantage. Elle est souvent plus lente et plus basse après le rebond.

Cela gêne l'adversaire car :
\begin{itemize}
    \item le rythme est cassé ;
    \item la balle rebondit bas ;
    \item l'adversaire doit plier davantage les jambes ;
    \item la balle peut être difficile à attaquer.
\end{itemize}

Le slice est donc un coup de variation et de perturbation.
\end{tennis}

\begin{oral}{rédaction possible}
Dans un slice, la rotation de la balle est différente de celle du lift. Le sens de la force de Magnus peut alors être orienté vers le haut. L'accélération verticale effective est réduite, ce qui peut donner à la balle une trajectoire plus flottante. Au tennis, cette balle ralentit le jeu et rebondit souvent plus bas, ce qui force l'adversaire à s'adapter.
\end{oral}

%=========================================================
\section{Piste 5 -- Rebond vertical : coefficient de restitution}

\subsection{Idée physique}

\begin{idee}{rebond imparfait}
Lors du rebond, la balle se déforme puis reprend sa forme. Une partie de l'énergie est restituée, mais une autre partie est perdue en chaleur, son et déformations internes.
\end{idee}

\subsection{Modèle}

\begin{modele}{chute verticale}
On lâche une balle sans vitesse initiale depuis une hauteur $h_0$.

Elle rebondit et atteint une hauteur $h_1$.

On néglige les frottements de l'air.
\end{modele}

\subsection{Bilan énergétique avant le choc}

Juste avant le choc, l'énergie potentielle initiale a été transformée en énergie cinétique :

\[
mgh_0=\frac12m(v^-)^2.
\]

Donc :

\[
v^-=\sqrt{2gh_0}.
\]

\subsection{Coefficient de restitution}

On définit le coefficient de restitution vertical :

\[
e=\frac{|v^+|}{|v^-|},
\]

où :
\begin{itemize}
    \item $v^-$ est la vitesse verticale juste avant le choc ;
    \item $v^+$ est la vitesse verticale juste après le choc.
\end{itemize}

Donc :

\[
v^+=ev^-.
\]

\subsection{Hauteur après rebond}

Après le rebond, la balle monte jusqu'à une hauteur $h_1$.

À cette hauteur, sa vitesse verticale est nulle.

Donc :

\[
\frac12m(v^+)^2=mgh_1.
\]

Ainsi :

\[
h_1=\frac{(v^+)^2}{2g}.
\]

Or :

\[
v^+=ev^-.
\]

Donc :

\[
h_1=\frac{e^2(v^-)^2}{2g}.
\]

Comme :

\[
(v^-)^2=2gh_0,
\]

on obtient :

\[
h_1=e^2h_0.
\]

Donc :

\[
e=\sqrt{\frac{h_1}{h_0}}.
\]

\subsection{Conclusion tennis}

\begin{tennis}{qualité du rebond}
Le coefficient de restitution permet de quantifier la qualité du rebond.

Si $e$ est grand :
\begin{itemize}
    \item la balle rebondit haut ;
    \item elle restitue bien l'énergie ;
    \item le jeu peut être plus rapide.
\end{itemize}

Si $e$ est petit :
\begin{itemize}
    \item la balle rebondit peu ;
    \item elle perd beaucoup d'énergie ;
    \item le jeu est ralenti.
\end{itemize}

Cela dépend de la balle, de sa pression interne, de sa température et de la surface.
\end{tennis}

\begin{oral}{rédaction possible}
Pour étudier le rebond, on peut utiliser le coefficient de restitution. En partant d'une hauteur $h_0$, la balle arrive au sol avec une vitesse $\sqrt{2gh_0}$. Après le rebond, elle repart avec une vitesse plus faible, égale à $e$ fois la vitesse initiale. On montre alors que la hauteur du rebond vérifie $h_1=e^2h_0$. Cette relation permet de mesurer expérimentalement la perte d'énergie lors du rebond.
\end{oral}

%=========================================================
\section{Piste 6 -- Rebond avec rotation : pourquoi le lift rebondit haut ?}

\subsection{Idée physique}

\begin{idee}{translation et rotation}
Une balle liftée possède une énergie cinétique de translation et une énergie cinétique de rotation. Pendant le rebond, le frottement avec le sol peut transformer une partie de la rotation en mouvement de translation.
\end{idee}

\subsection{Modèle}

\begin{modele}{balle en rotation}
On considère une balle de rayon $R$, de masse $m$, arrivant au sol avec :
\begin{itemize}
    \item une vitesse horizontale $v_x^-$ ;
    \item une vitesse verticale $v_y^-$ ;
    \item une vitesse angulaire $\omega^-$.
\end{itemize}

On note $I$ son moment d'inertie.

Pour simplifier :

\[
I=k mR^2,
\]

où $k$ dépend de la répartition de la masse.
\end{modele}

\subsection{Énergie avant le rebond}

L'énergie cinétique totale avant rebond est :

\[
E^-=\frac12m(v_x^-)^2+\frac12m(v_y^-)^2+\frac12I(\omega^-)^2.
\]

Après le rebond :

\[
E^+=\frac12m(v_x^+)^2+\frac12m(v_y^+)^2+\frac12I(\omega^+)^2.
\]

En réalité :

\[
E^+<E^-.
\]

Il y a perte d'énergie.

\subsection{Rôle du frottement}

Pendant le choc, le sol exerce :
\begin{itemize}
    \item une réaction normale $\vec{N}$ ;
    \item une force de frottement $\vec{f}$ tangentielle.
\end{itemize}

La force de frottement modifie :
\begin{itemize}
    \item la vitesse horizontale de la balle ;
    \item la vitesse angulaire de rotation.
\end{itemize}

\subsection{Lien translation--rotation}

Une relation importante est la condition de roulement sans glissement :

\[
v_x=R\omega.
\]

Si la balle glisse, cette relation n'est pas vérifiée.

Pour une balle liftée, la rotation vers l'avant peut augmenter l'effet horizontal après rebond.

\subsection{Conclusion tennis}

\begin{tennis}{rebond du lift}
Une balle liftée ne rebondit pas comme une balle plate.

À cause de sa rotation :
\begin{itemize}
    \item elle accroche davantage le sol ;
    \item elle peut repartir plus vite vers l'avant ;
    \item elle rebondit souvent plus haut ;
    \item elle force l'adversaire à frapper au-dessus de sa zone confortable.
\end{itemize}

Le lift est donc efficace non seulement dans l'air, mais aussi après le rebond.
\end{tennis}

\begin{oral}{rédaction possible}
Le rebond d'une balle liftée est plus complexe qu'un rebond vertical simple. La balle possède une énergie de translation, mais aussi une énergie de rotation. Lors du contact avec le sol, le frottement tangent modifie à la fois la vitesse horizontale et la rotation. Cela explique pourquoi une balle liftée peut repartir plus haut et plus vite après le rebond, ce qui rend le coup très difficile à contrôler pour l'adversaire.
\end{oral}

\begin{prepa}{théorème du moment cinétique}
En prépa, on peut utiliser le théorème du moment cinétique par rapport au centre de la balle :

\[
\sum \vec{M}_G(\vec{F}_{ext})=\frac{d\vec{L}_G}{dt}.
\]

La force de frottement du sol crée un moment qui modifie la vitesse angulaire $\omega$.
\end{prepa}

%=========================================================
\section{Piste 7 -- Influence du vent : les joueurs râlent-ils pour une bonne raison ?}

\subsection{Idée physique}

\begin{idee}{vitesse relative}
Les forces aérodynamiques dépendent de la vitesse de la balle par rapport à l'air, pas seulement de sa vitesse par rapport au sol. Le vent modifie donc la traînée et l'effet Magnus.
\end{idee}

\subsection{Modèle}

\begin{modele}{vitesse relative}
On note :
\begin{itemize}
    \item $\vec{v}$ : vitesse de la balle par rapport au sol ;
    \item $\vec{v}_{vent}$ : vitesse du vent par rapport au sol ;
    \item $\vec{v}_{rel}$ : vitesse de la balle par rapport à l'air.
\end{itemize}

Alors :

\[
\vec{v}_{rel}=\vec{v}-\vec{v}_{vent}.
\]
\end{modele}

\subsection{Bilan des forces}

\begin{bilan}{forces aérodynamiques avec vent}
La traînée devient :

\[
\vec{F}_D=-k\norm{\vec{v}_{rel}}\vec{v}_{rel}.
\]

Cette écriture est équivalente à une force en $v^2$, mais elle tient compte de la direction du vent.
\end{bilan}

\subsection{Cas 1 : vent de face}

Si la balle va vers la droite et que le vent souffle vers la gauche :

\[
\vec{v}=v\vec{i},
\qquad
\vec{v}_{vent}=-u\vec{i}.
\]

Alors :

\[
\vec{v}_{rel}=v\vec{i}-(-u\vec{i})=(v+u)\vec{i}.
\]

La norme est :

\[
v_{rel}=v+u.
\]

La force de traînée vaut en norme :

\[
F_D=k(v+u)^2.
\]

Comme :

\[
(v+u)^2>v^2,
\]

la traînée augmente.

\subsection{Cas 2 : vent dans le dos}

Si le vent souffle dans le même sens que la balle :

\[
\vec{v}_{vent}=u\vec{i}.
\]

Alors :

\[
\vec{v}_{rel}=(v-u)\vec{i}.
\]

La force de traînée vaut :

\[
F_D=k(v-u)^2.
\]

Elle est plus faible que sans vent si $u>0$.

\subsection{Cas 3 : vent latéral}

Si la balle avance selon $\vec{i}$ et que le vent souffle selon $\vec{j}$ :

\[
\vec{v}=v\vec{i},
\qquad
\vec{v}_{vent}=u\vec{j}.
\]

Alors :

\[
\vec{v}_{rel}=v\vec{i}-u\vec{j}.
\]

La force de traînée a une composante latérale, ce qui peut dévier la balle.

\subsection{Conclusion tennis}

\begin{tennis}{vent et précision}
Oui, les joueurs ont une raison physique de se plaindre du vent.

Le vent :
\begin{itemize}
    \item modifie la vitesse relative de la balle ;
    \item change l'intensité de la traînée ;
    \item peut amplifier ou réduire les effets ;
    \item peut dévier la balle latéralement ;
    \item gêne fortement le lancer de balle au service.
\end{itemize}

Comme le tennis se joue parfois à quelques centimètres près, une petite variation de trajectoire peut transformer un bon coup en faute.
\end{tennis}

\begin{oral}{rédaction possible}
Le vent intervient dans les forces aérodynamiques parce que celles-ci dépendent de la vitesse de la balle par rapport à l'air. Si le vent est de face, la vitesse relative augmente et la traînée devient plus forte. Si le vent est dans le dos, la traînée diminue et la balle peut sortir plus facilement. Si le vent est latéral, la force de traînée peut avoir une composante latérale. Cela justifie physiquement la gêne ressentie par les joueurs.
\end{oral}

%=========================================================
\section{Piste 8 -- Service : comment frapper fort tout en gardant la balle dans le carré ?}

\subsection{Idée physique}

\begin{idee}{contrainte géométrique}
Au service, le joueur doit frapper très fort, passer au-dessus du filet, puis faire retomber la balle dans un carré de service relativement court. Le lift ou le kick permet d'augmenter la sécurité.
\end{idee}

\subsection{Modèle}

\begin{modele}{service simplifié}
On modélise un service par une balle lancée depuis une hauteur $y_0$, avec une vitesse initiale $v_0$ et un angle $\alpha$.

On étudie d'abord la trajectoire sans effet :

\[
y(x)=y_0+x\tan(\alpha)-\frac{g}{2v_0^2\cos^2(\alpha)}x^2.
\]
\end{modele}

\subsection{Condition pour passer le filet}

Si le filet est à la distance $x_f$ et de hauteur $h_f$, il faut :

\[
y(x_f)>h_f.
\]

Donc :

\[
y_0+x_f\tan(\alpha)-\frac{g}{2v_0^2\cos^2(\alpha)}x_f^2>h_f.
\]

\subsection{Condition pour retomber dans le carré}

Si la limite du carré de service est à la distance $x_c$, il faut que la balle touche le sol avant ou à cette distance :

\[
y(x_c)\leq 0.
\]

Donc :

\[
y_0+x_c\tan(\alpha)-\frac{g}{2v_0^2\cos^2(\alpha)}x_c^2\leq 0.
\]

\subsection{Effet du lift au service}

Avec lift, on remplace $g$ par :

\[
g_{\text{eff}}=g+\frac{k_M}{m}v^2.
\]

La trajectoire devient approximativement :

\[
y(x)=y_0+x\tan(\alpha)-\frac{g_{\text{eff}}}{2v_0^2\cos^2(\alpha)}x^2.
\]

Comme :

\[
g_{\text{eff}}>g,
\]

le terme négatif est plus grand en valeur absolue.

Donc la balle descend plus vite.

\subsection{Conclusion tennis}

\begin{tennis}{service lifté}
Le lift permet au serveur de concilier puissance et sécurité.

Avec le lift :
\begin{itemize}
    \item la balle peut passer plus haut au-dessus du filet ;
    \item elle retombe plus vite dans le carré ;
    \item elle rebondit plus haut ;
    \item elle devient plus difficile à retourner.
\end{itemize}

C'est le principe du kick serve.
\end{tennis}

\begin{oral}{rédaction possible}
Le service est un problème de contraintes. La balle doit passer au-dessus du filet, mais retomber avant la limite du carré. Sans effet, une balle très rapide peut avoir une trajectoire trop tendue et sortir. Avec le lift, l'effet Magnus ajoute une force vers le bas. Dans le modèle, cela revient à remplacer $g$ par une pesanteur effective plus grande. La balle plonge donc plus vite, ce qui permet de frapper fort avec une marge de sécurité.
\end{oral}

%=========================================================
\section{Piste 9 -- Tension du cordage : puissance ou contrôle ?}

\subsection{Idée physique}

\begin{idee}{raquette comme ressort}
Lors de l'impact, les cordes de la raquette se déforment. Elles stockent temporairement de l'énergie élastique puis en restituent une partie à la balle. La tension du cordage modifie la raideur effective du système.
\end{idee}

\subsection{Modèle}

\begin{modele}{cordage assimilé à un ressort}
On assimile le cordage à un ressort de raideur $k$.

Lors de l'impact, la balle enfonce le cordage d'une distance maximale $x$.

La force de rappel exercée par le cordage est modélisée par la loi de Hooke :

\[
F=kx.
\]

L'énergie élastique stockée dans le cordage vaut :

\[
E_e=\frac12kx^2.
\]
\end{modele}

\subsection{Interprétation de la tension}

Un cordage plus tendu correspond à une raideur effective plus grande.

Donc :
\[
k \text{ grand} \quad \Rightarrow \quad \text{cordage rigide}.
\]

Un cordage moins tendu correspond à :
\[
k \text{ plus petit} \quad \Rightarrow \quad \text{cordage plus déformable}.
\]

\subsection{Mise en équation énergétique}

Avant l'impact, la balle arrive avec une énergie cinétique :

\[
E_c=\frac12mv^2.
\]

Pendant l'impact, une partie de cette énergie est stockée dans le cordage :

\[
\frac12mv^2 \approx \frac12kx^2 + E_{pertes}.
\]

Si on néglige temporairement les pertes :

\[
\frac12mv^2 \approx \frac12kx^2.
\]

Donc :

\[
x\approx v\sqrt{\frac{m}{k}}.
\]

\subsection{Conséquence physique}

Si $k$ augmente, alors :

\[
x\approx v\sqrt{\frac{m}{k}}
\]

diminue.

Donc un cordage très tendu se déforme moins.

Si $k$ diminue, le cordage se déforme davantage.

\subsection{Temps de contact}

Le système balle-cordage peut être comparé à un oscillateur masse-ressort.

La pulsation propre vaut :

\[
\omega_0=\sqrt{\frac{k}{m}}.
\]

La durée caractéristique de contact est de l'ordre :

\[
\tau \sim \frac{1}{\omega_0}=\sqrt{\frac{m}{k}}.
\]

Donc :
\begin{itemize}
    \item si $k$ est grand, le contact est court ;
    \item si $k$ est petit, le contact est plus long.
\end{itemize}

\subsection{Puissance et contrôle}

Un cordage moins tendu se déforme davantage et peut donner une sensation d'effet trampoline.
Il peut restituer davantage d'énergie à la balle, donc favoriser la puissance.

Un cordage très tendu se déforme moins.
La balle reste moins longtemps dans le cordage.
Le joueur ressent davantage de contrôle, mais moins de puissance gratuite.

\subsection{Lien avec les effets}

Pour donner du lift, il faut que les cordes accrochent la balle.

Le frottement entre cordes et balle transmet une force tangentielle qui crée un moment.

Si $f$ est la force tangentielle et $R$ le rayon de la balle, le moment par rapport au centre est :

\[
M=fR.
\]

Le théorème du moment cinétique donne qualitativement :

\[
M=I\frac{d\omega}{dt}.
\]

Donc :

\[
fR=I\frac{d\omega}{dt}.
\]

Plus la force tangentielle est importante, plus la variation de rotation est importante.

\subsection{Conclusion tennis}

\begin{tennis}{tension du cordage}
La tension du cordage influence fortement le jeu.

Un cordage peu tendu :
\begin{itemize}
    \item se déforme davantage ;
    \item donne plus de puissance ;
    \item augmente le temps de contact ;
    \item peut faciliter certains effets.
\end{itemize}

Un cordage très tendu :
\begin{itemize}
    \item se déforme moins ;
    \item donne plus de contrôle ;
    \item réduit l'effet trampoline ;
    \item demande plus d'engagement physique.
\end{itemize}

Le choix de tension est donc un compromis entre puissance, contrôle, confort et prise d'effet.
\end{tennis}

\begin{oral}{rédaction possible}
On peut modéliser le cordage d'une raquette comme un ressort. Lors de l'impact, la balle déforme les cordes, qui stockent une énergie élastique $\frac12kx^2$. Si le cordage est très tendu, la raideur $k$ est grande, donc la déformation est plus faible. À l'inverse, un cordage moins tendu se déforme davantage et donne un effet trampoline. Cela explique le compromis entre puissance et contrôle recherché par les joueurs.
\end{oral}

\begin{prepa}{modèle masse-ressort amorti}
Un modèle plus réaliste serait un oscillateur amorti :

\[
m\ddot{x}+c\dot{x}+kx=0.
\]

Le terme $c\dot{x}$ représente les pertes d'énergie pendant l'impact. Ce modèle permettrait d'étudier la restitution d'énergie et le confort de frappe.
\end{prepa}

%=========================================================
\section{Piste 10 -- Force de la raquette : impulsion et variation de quantité de mouvement}

\subsection{Idée physique}

\begin{idee}{frappe courte mais intense}
La frappe dure très peu de temps, mais la force exercée par la raquette peut être très grande. La bonne grandeur à étudier est alors l'impulsion.
\end{idee}

\subsection{Modèle}

\begin{modele}{impact raquette-balle}
On considère une balle arrivant vers la raquette avec une vitesse $\vec{v_i}$ et repartant avec une vitesse $\vec{v_f}$.

La durée de contact est notée $\Delta t$.
\end{modele}

\subsection{Quantité de mouvement}

La quantité de mouvement est :

\[
\vec{p}=m\vec{v}.
\]

La variation de quantité de mouvement est :

\[
\Delta \vec{p}=m\vec{v_f}-m\vec{v_i}.
\]

\subsection{Lien avec la force moyenne}

D'après la deuxième loi de Newton sous forme impulsionnelle :

\[
\vec{F}_{moy}\Delta t=\Delta \vec{p}.
\]

Donc :

\[
\vec{F}_{moy}=\frac{m(\vec{v_f}-\vec{v_i})}{\Delta t}.
\]

\subsection{Exemple numérique simple}

Si :
\[
m=\SI{0.058}{kg},
\]

et que la balle passe de :

\[
v_i=-\SI{20}{m.s^{-1}}
\]

à :

\[
v_f=\SI{40}{m.s^{-1}},
\]

alors :

\[
\Delta v=60\ \si{m.s^{-1}}.
\]

La variation de quantité de mouvement vaut :

\[
\Delta p=m\Delta v=0.058\times 60=3.48\ \si{kg.m.s^{-1}}.
\]

Si le contact dure :

\[
\Delta t=\SI{5e-3}{s},
\]

alors :

\[
F_{moy}=\frac{3.48}{5\times 10^{-3}}=696\ \si{N}.
\]

La force moyenne peut donc être très grande.

\subsection{Conclusion tennis}

\begin{tennis}{impact de la frappe}
La raquette exerce une force très intense pendant un temps très court.

Cela explique :
\begin{itemize}
    \item pourquoi la technique de frappe est importante ;
    \item pourquoi la rigidité de la raquette influence la sensation ;
    \item pourquoi le centrage de la balle est essentiel ;
    \item pourquoi les chocs répétés peuvent provoquer des douleurs.
\end{itemize}
\end{tennis}

\begin{oral}{rédaction possible}
Lors d'une frappe, il est difficile d'étudier la force instantanée car le contact est très court. On utilise donc l'impulsion. La variation de quantité de mouvement de la balle est égale au produit de la force moyenne par la durée de contact. Comme cette durée est de l'ordre de quelques millisecondes, la force moyenne peut être très grande, même pour une balle légère.
\end{oral}

%=========================================================
\section{Piste 11 -- Énergie d'une balle de tennis}

\subsection{Idée physique}

\begin{idee}{translation, rotation et pertes}
Une balle de tennis possède une énergie cinétique de translation, une énergie de rotation et une énergie potentielle de pesanteur. Ces énergies se transforment, mais une partie est dissipée dans l'air, le sol et la raquette.
\end{idee}

\subsection{Énergie de translation}

La balle de masse $m$ et de vitesse $v$ possède :

\[
E_c=\frac12mv^2.
\]

Exemple avec :

\[
m=\SI{0.058}{kg},
\qquad
v=\SI{50}{m.s^{-1}}.
\]

Alors :

\[
E_c=\frac12\times 0.058\times 50^2.
\]

\[
E_c=72.5\ \si{J}.
\]

\subsection{Énergie de rotation}

Si la balle tourne à la vitesse angulaire $\omega$ :

\[
E_{rot}=\frac12I\omega^2.
\]

Avec :

\[
I=kmR^2.
\]

Donc :

\[
E_{rot}=\frac12kmR^2\omega^2.
\]

\subsection{Énergie potentielle}

À une hauteur $h$ :

\[
E_p=mgh.
\]

\subsection{Bilan énergétique global}

Pendant un échange, on peut écrire qualitativement :

\[
E_{initiale}=E_{translation}+E_{rotation}+E_p+E_{pertes}.
\]

Les pertes viennent :
\begin{itemize}
    \item des frottements de l'air ;
    \item de la déformation de la balle ;
    \item du rebond ;
    \item du son ;
    \item de l'échauffement ;
    \item du cordage.
\end{itemize}

\subsection{Conclusion tennis}

\begin{tennis}{énergie et stratégie}
La puissance d'une frappe ne se résume pas à la vitesse.

Une partie de l'énergie peut être utilisée :
\begin{itemize}
    \item pour accélérer la balle ;
    \item pour la faire tourner ;
    \item pour produire du lift ;
    \item pour modifier le rebond.
\end{itemize}

Un joueur ne cherche donc pas toujours à maximiser la vitesse, mais à répartir intelligemment l'énergie entre translation et rotation.
\end{tennis}

\begin{oral}{rédaction possible}
L'énergie permet de comprendre qu'une frappe de tennis n'est pas seulement une question de vitesse. La balle possède une énergie cinétique de translation, mais aussi une énergie de rotation lorsqu'elle est liftée ou slicée. Une partie de l'énergie est également perdue dans les frottements et les rebonds. Le joueur doit donc répartir l'énergie de sa frappe entre puissance, effet et contrôle.
\end{oral}

%=========================================================
\section{Piste 12 -- Simulation numérique : résoudre quand l'équation devient trop difficile}

\subsection{Idée physique}

\begin{idee}{modéliser par ordinateur}
Lorsque l'on ajoute frottements, effet Magnus et vent, les équations deviennent difficiles à résoudre exactement. On peut alors les résoudre numériquement.
\end{idee}

\subsection{Modèle complet simplifié}

On écrit les forces :

\[
\vec{P}=m\vec{g},
\]

\[
\vec{F}_D=-k\norm{\vec{v}_{rel}}\vec{v}_{rel},
\]

\[
\vec{F}_M=k_M v^2 \vec{n}.
\]

La deuxième loi de Newton donne :

\[
m\vec{a}=m\vec{g}-k\norm{\vec{v}_{rel}}\vec{v}_{rel}+k_Mv^2\vec{n}.
\]

Donc :

\[
\vec{a}=\vec{g}-\frac{k}{m}\norm{\vec{v}_{rel}}\vec{v}_{rel}+\frac{k_M}{m}v^2\vec{n}.
\]

\subsection{Méthode d'Euler}

On choisit un pas de temps $\Delta t$.

À l'instant $t_n$ :
\[
\vec{v}_{n+1}=\vec{v}_n+\vec{a}_n\Delta t,
\]

\[
\vec{r}_{n+1}=\vec{r}_n+\vec{v}_n\Delta t.
\]

\subsection{Pseudo-code}

\begin{verbatim}
Initialiser position r et vitesse v
Choisir dt
Tant que y >= 0 :
    calculer v_rel = v - v_vent
    calculer le poids
    calculer la trainee
    calculer la force de Magnus
    calculer a = somme_forces / m
    v = v + a*dt
    r = r + v*dt
    enregistrer r
Tracer la trajectoire
\end{verbatim}

\subsection{Conclusion tennis}

\begin{tennis}{intérêt de la simulation}
La simulation permet de comparer :
\begin{itemize}
    \item balle sans frottements ;
    \item balle avec frottements ;
    \item balle liftée ;
    \item balle slicée ;
    \item balle avec vent.
\end{itemize}

C'est une méthode très pertinente pour un Grand Oral car elle montre la transition entre le modèle théorique et le tennis réel.
\end{tennis}

\begin{oral}{rédaction possible}
Dès que le modèle devient réaliste, les équations sont difficiles à résoudre à la main. On peut alors utiliser une méthode numérique comme la méthode d'Euler. L'idée est de découper le mouvement en petits intervalles de temps, puis de mettre à jour la vitesse et la position à chaque étape. Cela permet de comparer différentes trajectoires et de visualiser l'effet du lift ou du vent.
\end{oral}

%=========================================================
\section{Synthèse des pistes}

\begin{center}
\renewcommand{\arraystretch}{1.5}
\begin{tabular}{p{3.5cm}p{5cm}p{5.5cm}}
\toprule
\textbf{Piste} & \textbf{Équation principale} & \textbf{Impact tennis} \\
\midrule
Projectile sans frottements
&
$y(x)=y_0+x\tan\alpha-\dfrac{g}{2v_0^2\cos^2\alpha}x^2$
&
Modèle de base, trajectoire parabolique.
\\
\midrule
Frottements de l'air
&
$v(t)=\dfrac{v_0}{1+\frac{kv_0}{m}t}$
&
La balle ralentit fortement.
\\
\midrule
Lift
&
$a_y=-g-\dfrac{k_M}{m}v^2$
&
La balle plonge plus vite.
\\
\midrule
Slice
&
$a_y=-g+\dfrac{k_M}{m}v^2$
&
La balle flotte davantage et rebondit bas.
\\
\midrule
Rebond vertical
&
$h_1=e^2h_0$
&
Mesure de la perte d'énergie.
\\
\midrule
Rebond avec rotation
&
$E=\dfrac12mv^2+\dfrac12I\omega^2$
&
Rotation et rebond sont couplés.
\\
\midrule
Vent
&
$\vec{v}_{rel}=\vec{v}-\vec{v}_{vent}$
&
La trajectoire devient moins prévisible.
\\
\midrule
Cordage
&
$E_e=\dfrac12kx^2$
&
Compromis puissance--contrôle.
\\
\midrule
Frappe
&
$\vec{F}_{moy}=\dfrac{m(\vec{v_f}-\vec{v_i})}{\Delta t}$
&
Force très grande sur temps très court.
\\
\bottomrule
\end{tabular}
\end{center}

%=========================================================
\section{Conclusion générale}

L'étude d'une balle de tennis permet de relier des notions de Terminale à des notions plus avancées de prépa.

Le modèle de base est celui du projectile soumis à son poids. Il donne une trajectoire parabolique. Mais pour comprendre réellement le tennis, il faut enrichir ce modèle avec :
\begin{itemize}
    \item les frottements de l'air ;
    \item la rotation ;
    \item l'effet Magnus ;
    \item le rebond ;
    \item le vent ;
    \item la déformation du cordage ;
    \item les transferts d'énergie.
\end{itemize}

\begin{oral}{conclusion possible pour le Grand Oral}
Finalement, la physique montre que le tennis n'est pas seulement un sport de puissance. Un joueur ne cherche pas seulement à frapper fort : il doit contrôler la trajectoire, la rotation, le rebond et la profondeur de balle. Le lift, le slice, le choix du cordage ou l'adaptation au vent sont autant de stratégies qui reposent sur des phénomènes physiques. Une simple balle de tennis devient donc un objet d'étude très riche, qui relie directement la mécanique de Terminale à des notions de prépa comme les équations différentielles, l'énergie de rotation, l'impulsion et la dynamique des fluides.
\end{oral}

%=========================================================
\section{Mini-fiche de formules}

\begin{center}
\renewcommand{\arraystretch}{1.6}
\begin{tabular}{p{5cm}p{8cm}}
\toprule
\textbf{Notion} & \textbf{Formule} \\
\midrule
Deuxième loi de Newton & $\sum \vec{F}=m\vec{a}$ \\
Poids & $\vec{P}=m\vec{g}$ \\
Trajectoire sans frottements & $y(x)=y_0+x\tan\alpha-\dfrac{g}{2v_0^2\cos^2\alpha}x^2$ \\
Traînée & $\vec{F}_D=-k\norm{\vec{v}}\vec{v}$ \\
Vitesse avec frottement quadratique & $v(t)=\dfrac{v_0}{1+\frac{kv_0}{m}t}$ \\
Force de Magnus & $F_M=k_Mv^2$ \\
Pesanteur effective avec lift & $g_{\text{eff}}=g+\dfrac{k_M}{m}v^2$ \\
Coefficient de restitution & $e=\dfrac{|v^+|}{|v^-|}$ \\
Hauteur de rebond & $h_1=e^2h_0$ \\
Énergie cinétique & $E_c=\dfrac12mv^2$ \\
Énergie de rotation & $E_{rot}=\dfrac12I\omega^2$ \\
Énergie élastique du cordage & $E_e=\dfrac12kx^2$ \\
Impulsion & $\vec{F}_{moy}\Delta t=\Delta \vec{p}$ \\
Vitesse relative au vent & $\vec{v}_{rel}=\vec{v}-\vec{v}_{vent}$ \\
Méthode d'Euler & $\vec{v}_{n+1}=\vec{v}_n+\vec{a}_n\Delta t$ \\
\bottomrule
\end{tabular}
\end{center}

\end{document}